\documentclass[12pt]{article}
\usepackage[utf8]{inputenc}
\usepackage[T1]{fontenc}
\usepackage[margin=1in]{geometry}
\usepackage{setspace}
\usepackage{xcolor}
\usepackage[colorlinks=true, urlcolor=blue!60!black, linkcolor=blue!60!black]{hyperref}
\setlength{\parindent}{0em}
\setlength{\parskip}{0.6em}
\renewcommand{\baselinestretch}{1.0}
\begin{document}
\begin{center}
{\LARGE \textit{Students for Fair Admissions, Inc. v. President and Fellows of Harvard College}, 600 U.S. 181 (2023)\par}
\vspace{0.5em}
{\large Case Brief --- Two-Page Edition\par}
\end{center}

\section{Memory Jogger}
Race-conscious admissions at Harvard and UNC violate equal protection and Title VI, though applicants may discuss how race affected their lives (\href{https://www.law.cornell.edu/supremecourt/text/20-1199}{600 U.S.\ 181}).

\section{Detailed Case Facts}
Students for Fair Admissions sued Harvard, a private university subject to Title VI, and UNC, a public university. Both considered race in holistic review: Harvard tracked the racial makeup of its tentative class and showed race on its final ``lop'' list, and UNC readers could award a racial ``plus.'' Both relied on the diversity interest recognized in \textit{Bakke} and \textit{Grutter}.

\section{Procedural History}
After bench trials, the district courts upheld both programs; the First Circuit affirmed in \textit{Harvard}. The Court granted review in both.

\section{Judicial Votes}
Reversed, 6--2 (\textit{Harvard}) and 6--3 (\textit{UNC}). Roberts wrote for Thomas, Alito, Gorsuch, Kavanaugh, and Barrett; Thomas, Gorsuch, and Kavanaugh concurred; Sotomayor and Jackson dissented; Jackson was recused in \textit{Harvard}.

\section{Holding}
Both programs fail strict scrutiny: their goals are unmeasurable, they use race as a negative and a stereotype, and they lack a logical endpoint. Universities may consider an applicant's discussion of race tied to his or her character or ability. Military academies were reserved.

\section{Analysis of Opinions}
Roberts held that the Equal Protection Clause is universal and found the programs' goals ``commendable'' but too elusive for review, their racial categories imprecise, their zero-sum preferences a minus for others, and their stable racial shares proof of no endpoint. Thomas defended a colorblind original meaning; Gorsuch read Title VI's text to bar any racial disadvantage; Kavanaugh relied on \textit{Grutter}'s 25-year horizon. Sotomayor accused the Court of ``cement[ing] a superficial rule of colorblindness'' and invoked race-conscious Reconstruction laws; Jackson documented ``gulf-sized race-based gaps.''

\section{Examples: Future Applications}
\textbf{Same-Side:} holding minority enrollment in a fixed band at a medical school is invalid. \textbf{Opposite-Side:} crediting an essay about leadership after facing discrimination is valid. \textbf{Fence-Sitter:} race-conscious admissions at military academies.

\section{Critique}
Originalists praise the anti-classification principle but note the unexplained military carve-out. Progressive critics invoke an anti-subordination reading of Reconstruction history and fault the Court for ignoring trial findings and reliance interests.

\section{Key Quotations}
\begin{quote}
``Eliminating racial discrimination means eliminating all of it. And the Equal Protection Clause, we have accordingly held, applies `without regard to any differences of race, of color, or of nationality'---it is `universal in [its] application.' ''
\end{quote}
\noindent --- Chief Justice Roberts, Opinion of the Court, 600 U.S.\ at 206.

\end{document}
